\section{Fishbowl Assessment Procedure}

This section outlines the fishbowl assessment so that you know what will happen, why it matters, and how to prepare. In the fishbowl, one group discusses its project in the centre of the room while other students observe, question, and respond from the outside.

Each group will normally receive about 30 minutes to explain how it applied innovation and product development principles while developing the selected hydroponics-related product. The goal is to assess more than speaking confidence. The discussion should show that you understand the decisions your group made, the evidence that supports those decisions, and the trade-offs that still remain.

\subsection{How the Session Runs}
\begin{enumerate}
    \item Your group explains the project context, current concept, and key development decisions.
    \item Observing students ask constructive questions and challenge unclear reasoning.
    \item The presenting group responds, defends its choices, and acknowledges areas that still require improvement.
    \item The lecturer and peers provide feedback after the discussion.
\end{enumerate}

This process is intended to improve the project, not merely to criticise it. The session also gives you practice communicating technical ideas to an audience that may not share the same assumptions as your team.

\subsection{Assessment Criteria}

Your individual performance during the fishbowl will be judged using the following criteria:
\begin{itemize}
    \item \textbf{Body language and attentiveness}: you remain engaged and responsive throughout the discussion.
    \item \textbf{Confidence and clarity}: you explain ideas clearly and contribute actively.
    \item \textbf{Defence of ideas}: you justify decisions with logic, evidence, or project data.
    \item \textbf{Connection to group objectives}: your comments support the overall project direction rather than isolated opinions.
    \item \textbf{Critical thinking}: you respond thoughtfully to questions and recognise limitations or alternatives.
\end{itemize}

\subsection{Preparation Tips}
\begin{itemize}
    \item review your team's current needs, specifications, and concept choices before the session;
    \item prepare short explanations for important technical decisions;
    \item bring notes or supporting visuals if the lecturer allows them; and
    \item listen carefully to questions so that your answers remain relevant and precise.
\end{itemize}

\subsection{Reflection Prompt}

After the fishbowl, note which questions exposed weak reasoning or missing evidence in your project. Use that feedback to strengthen the final presentation and report.
